%% kapa: exam=1 round=22 year=2011 period=1 form=A subject=civil_law subject_ko=민법 seq=23 number=23
%% source: Q-Net 감정평가사 기출문제 (cst00309) artlSeq=5036790 file=extracted/1008/5036790/2011년 제22회 감정평가사 1차시험 기출문제/1교시.hwp
%% conversion: kordoc 4.24.51 -> tools/gen1.py (auto); review=proofread-sonnet; flags=none
\begin{question}{23}
甲의 아들 乙은 무단으로 甲의 대리인이라고 칭하며 甲소유의 토지를 丙에게 매도하였다. 다음 설명으로 옳은 것은?
\begin{choices}
  \item 丙은 甲이 추인한 후라도 매매계약을 철회할 수 있다.
  \item 甲이 乙에게 추인의 의사표시를 한 경우, 丙이 그 사실을 모르더라도 甲의 매매대금지급청구에 대해 丙은 대항할 수 없다.
  \item 甲이 매매계약의 추인을 거절한 후에 사망한 경우, 乙이 甲을 단독 상속하였다고 해도 무권대리행위가 유효하게 되지는 않는다.
  \item 乙이 미성년자인 경우 甲이 매매계약의 추인을 거절한다면, 丙은 乙에게 계약이행을 청구할 수 있으나 손해배상을 청구할 수는 없다.
  \item 丙이 상당한 기간을 정하여 매매계약의 추인 여부를 최고하였으나 甲이 그 기간 내에 확답을 발하지 않으면 추인한 것으로 본다.
\end{choices}
\end{question}
\answer{3}
